\documentclass[11pt]{amsart}

\usepackage[T1]{fontenc}
\usepackage{lmodern}
\usepackage{microtype}
\usepackage{mathtools,amssymb,amsthm}
\usepackage[hidelinks]{hyperref}

\hypersetup{
  pdftitle={The Holtz--Sturmfels Ideal-Generation Conjecture Fails for n at Least 5},
  pdfkeywords={principal minors, symmetric matrices, hyperdeterminant, invariant theory}
}

\newtheorem{theorem}{Theorem}
\newtheorem{proposition}[theorem]{Proposition}
\newtheorem{lemma}[theorem]{Lemma}
\newtheorem{corollary}[theorem]{Corollary}
\theoremstyle{remark}
\newtheorem{remark}[theorem]{Remark}

\newcommand{\C}{\mathbb C}
\newcommand{\PP}{\mathbb P}
\newcommand{\HD}{\mathrm{HD}}
\newcommand{\Sym}{\operatorname{Sym}}
\newcommand{\bS}{\mathbf S}
\newcommand{\Span}{\operatorname{span}}

\title[Failure of the Holtz--Sturmfels conjecture]
{The Holtz--Sturmfels Ideal-Generation Conjecture Fails for \(n\geq5\)}

\date{}

\subjclass[2020]{13A50, 14M12, 15A15, 15A69}
\keywords{principal minors, symmetric matrix, hyperdeterminant, invariant theory}

\begin{document}

\begin{abstract}
Let \(P_n\) be the homogeneous prime ideal of the projective variety of
principal minors of symmetric \(n\times n\) matrices, and let \(J_n\) be the
ideal generated by the hyperdeterminantal module \(\HD(n)\).
Holtz and Sturmfels conjectured that \(J_n=P_n\).
We work over \(\C\).
For \(n=5\), we construct five \(\mathrm{SL}_2^5\)-invariant quartics
\(B_1,\dots,B_5\). Their pullbacks under the principal-minor map are all
\(12t^{10}\Pi^2\), where \(\Pi\) is the pentad. Hence the four differences
\(B_1-B_i\), \(2\leq i\leq5\), lie in \((P_5)_4\); a \(4\times4\)
coefficient determinant equal to \(-512\) shows that they are linearly
independent. The module \(\HD(5)\) has dimension \(250\) and contains no
nonzero \(\mathrm{SL}_2^5\)-invariant, so
\(\dim(P_5)_4\geq254\) and \(J_5\neq P_5\).
The same relations on any five principal indices show that
\(J_n\neq P_n\) for every \(n\geq5\). Since Oeding proved
\(\sqrt{J_n}=P_n\), the hyperdeterminantal ideal is nonradical for every
\(n\geq5\). All polynomials exhibited have integer coefficients, and the
real form of the conjecture --- the form in which it was posed --- fails
for every \(n\geq5\) as well.
\end{abstract}

\maketitle

\section{Statement of the result}\label{sec:statement}

We work over \(\C\) throughout, as \cite{Oeding} does; \cite{HS} pose
Conjecture~14 over \(\mathbb R\), and the Status paragraph at the end
records the descent to that form.

For \(V_i=\C^2\), set
\[
  V^{(n)}=V_1\otimes\cdots\otimes V_n,
  \qquad
  R_n=\Sym\bigl((V^{(n)})^*\bigr)
      =\C[A_I:I\subseteq[n]].
\]
Let \(X_I\) denote the principal submatrix of a symmetric matrix \(X\)
indexed by \(I\). The homogeneous principal-minor map is
\[
  \phi_n:\PP\bigl(\Sym^2\C^n\oplus\C\bigr)
  \dashrightarrow \PP(V^{(n)}),
  \qquad
  [X:t]\longmapsto
  \bigl[t^{\,n-|I|}\det X_I\bigr]_{I\subseteq[n]}.
\]
Let \(Z_n\) be the closure of its image and let \(P_n=I(Z_n)\).
Let
\[
  G_n=\mathrm{SL}(V_1)\times\cdots\times\mathrm{SL}(V_n)
\]
act on \(R_n\) through \(V^{(n)}\), and let the symmetric group \(S_n\)
act on \(R_n\) by permuting the \(n\) tensor factors, hence the indices of
the coordinates \(A_I\). Write \(\HD(n)\) for the \emph{hyperdeterminantal
module}: the linear span of the orbit of Cayley's \(2\times2\times2\)
hyperdeterminant under \(G_n\rtimes S_n\). This is the module of
\cite[\S6]{HS} and the module \(\HD\) of \cite[\S5]{Oeding}. Equivalently,
\(\HD(n)\) is the \(G_n\)-module generated by the hyperdeterminants of all
\(\binom n32^{\,n-3}\) two-by-two-by-two subtensors of
\((A_I)_{I\subseteq[n]}\): by \eqref{eq:HDdecomp} and its analogue for
general \(n\) in the proof of Theorem~\ref{thm:main}, the summand attached
to a triple of tensor directions is an irreducible \(G_n\)-module and
contains the hyperdeterminants of all \(2^{\,n-3}\) subtensors that use
those three directions, while \(S_n\) permutes the \(\binom n3\) triples.
Both groups are needed: the \(G_n\)-orbit span of a single subtensor
hyperdeterminant is one summand of \eqref{eq:HDdecomp}, of dimension
\(25\) rather than \(250\) when \(n=5\). One has \(\HD(n)\subseteq(P_n)_4\)
\cite[Proposition~5.5]{Oeding}, and we set
\[
  J_n=(\HD(n))\subseteq R_n.
\]
Holtz and Sturmfels conjectured that \(J_n=P_n\)
\cite[Conjecture~14]{HS}. Oeding proved the set-theoretic equality
\(\mathcal V(J_n)=Z_n\) \cite[Theorem~1.3]{Oeding}; the ideal-theoretic
problem was subsequently recorded as open for \(n\geq5\)
\cite[p.~1194]{HO}.

We use Oeding's representation-theoretic description of \(\HD(n)\). In
particular,
\begin{equation}\label{eq:dimHD}
  \dim\HD(n)=\binom n3 5^{n-3}
  \qquad\text{\cite[Observation~5.1]{Oeding}.}
\end{equation}
The formula printed for the dimension of the same module in \cite[\S6]{HS}
is \(\binom{2^{\,n-3}+3}{4}\binom n3\). The two agree for \(n\leq4\),
giving \(1\) and \(20\), and disagree from \(n=5\) on, where they give
\(250\) and \(350\). The disagreement is one of counts and not of modules:
both papers define \(\HD(n)\) as the span of the orbit of the
\(2\times2\times2\) hyperdeterminant under \(G_n\) and \(S_n\), and
Conjecture~14 is the assertion that this span generates \(P_n\), so that
span is the object of Theorem~\ref{thm:main} below. Neither value is used
in the refutation. What separates our quartics from \(\HD(n)\) in
Proposition~\ref{prop:n5} and in the proof of Theorem~\ref{thm:main} is the
\(G_n\)-isotypic type of \(\HD(n)\), which forces \(\HD(n)^{G_n}=0\)
whatever the dimension of \(\HD(n)\) may be; \eqref{eq:dimHD} enters only
in the numerical bound printed in Theorem~\ref{thm:main}.

\begin{theorem}\label{thm:main}
For every \(n\geq5\),
\[
  J_n\subsetneq P_n.
\]
More precisely, \((P_5)_4\) contains a four-dimensional
\(\mathrm{SL}_2^5\)-trivial subspace \(W\) with \(W\cap\HD(5)=0\), and hence
\[
  \dim(P_5)_4\geq254>250=\dim\HD(5).
\]
Consequently, \(J_n\) is nonradical for every \(n\geq5\).
\end{theorem}

\section{Five invariant quartics}

Identify \(I\subseteq[5]\) with its indicator
\(\alpha\in\{0,1\}^5\), and write \(A_\alpha=A_I\).
Let \(\varepsilon\) be the alternating form on \(\C^2\), normalized by
\[
  \varepsilon_{01}=1,
  \qquad
  \varepsilon_{10}=-1,
  \qquad
  \varepsilon_{00}=\varepsilon_{11}=0.
\]
For \(i\in[5]\), define
\begin{equation}\label{eq:Qi}
  q^{(i)}_{ab}(A)
  =
  \sum_{\substack{\alpha,\beta\in\{0,1\}^5\\
                   \alpha_i=a,\ \beta_i=b}}
  A_\alpha A_\beta
  \prod_{j\neq i}\varepsilon_{\alpha_j\beta_j}.
\end{equation}
Set
\[
  Q_i(A)=\bigl(q^{(i)}_{ab}(A)\bigr)_{a,b=0}^1,
  \qquad
  B_i(A)=\det Q_i(A).
\]
The four contractions make \(Q_i(A)\) symmetric. This is the standard
construction of the five degree-four invariants of a five-qubit tensor;
see \cite{LT}.

\begin{lemma}\label{lem:invariance}
Each \(B_i\) is invariant under \(G_5\).
\end{lemma}

\begin{proof}
The alternating form is \(\mathrm{SL}_2\)-invariant. Contracting the
four factors other than \(i\) therefore gives
\[
  Q_i(g\cdot A)=g_iQ_i(A)g_i^{\mathsf T}
  \qquad(g=(g_1,\dots,g_5)\in G_5),
\]
where \(g_i\) acts through the identification of \(V_i\) with \(V_i^*\)
furnished by \(\varepsilon\); without that identification the right-hand
side reads \(g_i^{-\mathsf T}Q_i(A)g_i^{-1}\). Taking determinants and
using \(\det g_i=1\), which disposes of either form, proves the claim.
\end{proof}

\section{The quartic relations}\label{sec:relations}

For a symmetric \(5\times5\) matrix \(X=(x_{ij})\), define the pentad
\begin{align*}
\Pi(X)={}&
 x_{12}x_{13}x_{24}x_{35}x_{45}
-x_{12}x_{13}x_{25}x_{34}x_{45}
-x_{12}x_{14}x_{23}x_{35}x_{45}\\
&+x_{12}x_{14}x_{25}x_{34}x_{35}
+x_{12}x_{15}x_{23}x_{34}x_{45}
-x_{12}x_{15}x_{24}x_{34}x_{35}\\
&+x_{13}x_{14}x_{23}x_{25}x_{45}
-x_{13}x_{14}x_{24}x_{25}x_{35}
-x_{13}x_{15}x_{23}x_{24}x_{45}\\
&+x_{13}x_{15}x_{24}x_{25}x_{34}
+x_{14}x_{15}x_{23}x_{24}x_{35}
-x_{14}x_{15}x_{23}x_{25}x_{34}.
\end{align*}

\begin{proposition}\label{prop:pullback}
For every \(i\in[5]\),
\begin{equation}\label{eq:pullback}
  B_i\bigl(\phi_5(X,t)\bigr)=12t^{10}\Pi(X)^2
\end{equation}
as an identity in \(\mathbb Z[x_{ij},t:1\leq i\leq j\leq5]\).
\end{proposition}

\begin{proof}
If a summand in \eqref{eq:Qi} is nonzero, then
\(\alpha_j+\beta_j=1\) for every \(j\neq i\). Thus
\(q^{(i)}_{ab}\) has subset weight \(4+a+b\), and every monomial of
\(B_i\) has subset weight \(10\). It is therefore enough to verify
\eqref{eq:pullback} on the affine chart \(t=1\).

Substituting the thirty-two principal minors \(A_I=\det X_I\) into
\eqref{eq:Qi} and expanding over \(\mathbb Z\) gives
\[
  B_i\bigl((\det X_I)_{I\subseteq[5]}\bigr)=12\Pi(X)^2
  \qquad(i=1,\dots,5).
\]
This is an exact polynomial calculation, and it is specified completely by
what precedes: each side has exactly seventy-three monomials, all of degree
ten in the fifteen entries \(x_{jk}\), and each of the five differences
\(B_i\bigl((\det X_I)_{I\subseteq[5]}\bigr)-12\Pi(X)^2\) vanishes
coefficientwise. No step of the expansion is cancellation-sensitive, and it
is reproduced in exact integer arithmetic, with no sampling and no floating
point, by any computer algebra system.
\end{proof}

Set
\[
  F_i=B_1-B_i,
  \qquad i=2,3,4,5.
\]
Proposition~\ref{prop:pullback} implies
\begin{equation}\label{eq:membership}
  F_2,F_3,F_4,F_5\in(P_5)_4.
\end{equation}

\begin{lemma}\label{lem:independence}
The quartics \(F_2,F_3,F_4,F_5\) are linearly independent.
\end{lemma}

\begin{proof}
Consider the coefficients of
\[
\begin{aligned}
 m_1&=A_{134}A_{234}A_{15}A_{25},
&\qquad
 m_2&=A_{124}A_{234}A_{15}A_{35},\\
 m_3&=A_{124}A_{134}A_{25}A_{35},
&
 m_4&=A_{123}A_{234}A_{15}A_{45}.
\end{aligned}
\]
With columns ordered as \((F_2,F_3,F_4,F_5)\), the coefficient matrix is
\[
  \begin{pmatrix}
   0& 4& 4& 4\\
   4& 0& 4& 4\\
  -4&-4& 0& 0\\
   4& 4& 0& 4
  \end{pmatrix},
  \qquad
  \det=-512.
\]
Hence the four quartics are independent.
\end{proof}

\section{Separation from the hyperdeterminantal module}\label{sec:separation}

As a \(G_5\)-module, \(\HD(5)\) decomposes as
\begin{equation}\label{eq:HDdecomp}
 \HD(5)
 \cong
 \bigoplus_{\substack{T\subseteq[5]\\ |T|=3}}
 \left(
   \bigotimes_{j\in T}\bS_{(2,2)}V_j^*
 \right)
 \otimes
 \left(
   \bigotimes_{j\notin T}\bS_{(4)}V_j^*
 \right).
\end{equation}
This is Oeding's compact module notation, defined at \cite[p.~85]{Oeding}
as the sum over those permutations of a partition tuple that yield distinct
\(G_n\)-modules, and restated inside the proof of
\cite[Proposition~5.2]{Oeding}; that proposition shows in addition that
\(\HD\) occurs with multiplicity one in \(\Sym^4\bigl((V^{(n)})^*\bigr)\)
and is irreducible as a \(G_n\rtimes S_n\)-module. The dimension
\eqref{eq:dimHD} is the count of exactly this direct sum
\cite[Observation~5.1]{Oeding}.

Now \(\bS_{(2,2)}\C^2\) is one-dimensional and trivial on
\(\mathrm{SL}_2\), whereas
\(\bS_{(4)}\C^2\cong\Sym^4(\C^2)\) has dimension five and no nonzero
\(\mathrm{SL}_2\)-fixed vector. Therefore
\begin{equation}\label{eq:HDfacts}
  \dim\HD(5)=\binom53 5^2=250,
  \qquad
  \HD(5)^{G_5}=0.
\end{equation}

\begin{proposition}\label{prop:n5}
The degree-four component \((P_5)_4\) strictly contains \(\HD(5)\).
Indeed,
\[
  W=\Span\{F_2,F_3,F_4,F_5\}
\]
is a four-dimensional \(G_5\)-trivial subspace of \((P_5)_4\) and
\(W\cap\HD(5)=0\).
\end{proposition}

\begin{proof}
Membership and \(\dim W=4\) follow from
\eqref{eq:membership} and Lemma~\ref{lem:independence}.
Lemma~\ref{lem:invariance} shows that \(G_5\) fixes \(W\) pointwise,
while \eqref{eq:HDfacts} gives \(W\cap\HD(5)=0\). Since
\(\HD(5)\subseteq(P_5)_4\) by \cite[Proposition~5.5]{Oeding},
\[
  \dim(P_5)_4\geq\dim\HD(5)+\dim W=254.
\]
Finally, \((J_5)_4=\HD(5)\), because \(J_5\) is generated in degree
four. Hence \(J_5\neq P_5\).
\end{proof}

\begin{proof}[Proof of Theorem~\ref{thm:main}]
Proposition~\ref{prop:n5} proves the assertion for \(n=5\).
Let \(n>5\), fix a five-element set \(S\subseteq[n]\), and identify the
coordinates \(A_I\), \(I\subseteq S\), with the coordinates of
\(R_5\). Let \(F_i^S\in R_n\) be the resulting copy of \(F_i\).
For every symmetric \(n\times n\) matrix \(X\), the subvector of
\(\phi_n(X,t)\) indexed by \(I\subseteq S\) is
\(t^{\,n-5}\phi_5(X_S,t)\). Since \(F_i\) is homogeneous of degree four,
\[
  F_i^S\bigl(\phi_n(X,t)\bigr)
  =t^{4(n-5)}F_i\bigl(\phi_5(X_S,t)\bigr)=0.
\]
Thus \(F_i^S\in(P_n)_4\).

Set \(H_S=\prod_{j\in S}\mathrm{SL}(V_j)\). Each \(F_i^S\) is fixed by
\(H_S\). On the other hand, restriction of Oeding's decomposition gives
\[
 \HD(n)\big|_{\prod_{j=1}^n\mathrm{SL}(V_j)}
 \cong
 \bigoplus_{\substack{T\subseteq[n]\\ |T|=3}}
 \left(
   \bigotimes_{j\in T}\bS_{(2,2)}V_j^*
 \right)
 \otimes
 \left(
   \bigotimes_{j\notin T}\bS_{(4)}V_j^*
 \right).
\]
For every \(T\) with \(|T|=3\), at least two indices of \(S\) lie
outside \(T\). Thus every summand contains a nontrivial
\(\Sym^4(V_j^*)\) factor for the action of \(H_S\), and consequently
\[
  \HD(n)^{H_S}=0.
\]
Since \(F_i^S\neq0\), it follows that \(F_i^S\notin\HD(n)=(J_n)_4\).
Therefore \(J_n\neq P_n\) for every \(n\geq5\).

Finally, \(\mathcal V(J_n)=Z_n\) \cite[Theorem~1.3]{Oeding} and the
Nullstellensatz give \(\sqrt{J_n}=I(Z_n)=P_n\). Hence \(J_n\) is
nonradical for every \(n\geq5\).
\end{proof}

\medskip
\noindent\textbf{Exact verification.}
Every check described here is a recomputation from the data printed in
Sections~\ref{sec:statement}--\ref{sec:separation}, in exact integer
arithmetic and with no floating point, so a referee can redo it from this
paper alone. The computation has been re-run once, independently, on a
separate machine, by a program written from those sections alone, and all
\(36\) of its checks agree with what is reported here. It re-derives the
pentad as the signed sum over the twelve Hamiltonian cycles of \(K_5\) and
recovers the twelve monomials printed in Section~\ref{sec:relations}; it
verifies \eqref{eq:pullback} coefficientwise in
\(\mathbb Z[x_{jk},t]\) for all five \(i\), seventy-three monomials on each
side and no nonzero residual; it recomputes the matrix of
Lemma~\ref{lem:independence}, its determinant \(-512\), and the rank \(4\)
of the full \(4\times280\) coefficient matrix of \(F_2,\dots,F_5\); it
obtains \(\dim\bS_{(2,2)}\C^2=1\), \(\dim\bS_{(4)}\C^2=5\) and
\(\Sym^m(\C^2)^{\mathrm{SL}_2}=0\) for \(m>0\), hence
\eqref{eq:HDfacts}; it builds the torus-weight-zero part of \(\HD(5)\)
explicitly, finds its rank to be \(10\) and finds that adjoining
\(F_2,\dots,F_5\) raises the rank to \(14\), which confirms
\(W\cap\HD(5)=0\) by computation as well as by \eqref{eq:HDfacts} --- a
confirmation only, since the matching upper bound
\(\dim\HD(5)^{\mathrm{wt}\,0}\leq10\) was there obtained by sampling the
orbit, whereas the proof of Proposition~\ref{prop:n5} uses
\eqref{eq:HDfacts} and not this computation; it checks that Cayley's
hyperdeterminant of each of the forty \(2\times2\times2\) subtensors of
\(\phi_5(X,t)\) is the zero polynomial; and it re-runs the reduction in the
proof of Theorem~\ref{thm:main} for \(n=6,\dots,12\) over every one of the
\(1715\) five-element subsets and every \(i\), together with the module
facts for \(n=3,\dots,14\) over every pair of a three-element and a
five-element subset.

Two limits of that re-run should be recorded. Theorem~\ref{thm:main}
asserts \(J_n\neq P_n\) for every \(n\geq5\), an infinite family: the
re-run covers \(n=5\) exactly and symbolically, the relations for
\(n=6,\dots,12\) at exact integer points, and the module facts for
\(n\leq14\), but not the relations for \(n=13,14\) and nothing at all for
\(n>14\); those rest on the uniform argument in the proof of
Theorem~\ref{thm:main} and not on computation. And the inclusion
\(\HD(n)\subseteq(P_n)_4\) was re-derived only at the orbit generators for
\(n=5\), by the vanishing check just described, while \(\sqrt{J_n}=P_n\)
was taken as cited and not re-derived. The program and its transcript are
retained in an auxiliary archive and are available on request. The claim
recorded here is that an independent re-run agrees with what is reported
above, not that it constitutes a second proof.

\medskip
\noindent\textbf{Status and scope.}
We have worked over \(\C\), as \cite{Oeding} does, whereas Conjecture~14 is
posed in \cite{HS} over \(\mathbb R\), with ring \(\mathbb R[A_\bullet]\)
and group \(\mathrm{SL}_2(\mathbb R)^n\). The conclusion descends to that
form. The quartics \(F_i\) and \(F_i^S\) have integer coefficients and
\eqref{eq:pullback} is an identity over \(\mathbb Z\), so each is a
relation among the principal minors of a real symmetric matrix. And
\(\HD(n)\), the span of the orbit of a quartic with integer coefficients
under a group defined over \(\mathbb Q\), has a basis of quartics with
integer coefficients, so a nonzero real element of the intersection of
\(\HD(n)\) with the real span of the \(F_i^S\) would be a nonzero element
of the complex intersection, which is \(0\). Hence the real form of
Conjecture~14 fails for every \(n\geq5\) as well.

What is refuted is the ideal-theoretic statement, and only for
\(n\geq5\). Conjecture~14 holds for \(n=3\), where \(P_3\) is principal,
and for \(n=4\), where it is verified by the computations of
\cite[\S4]{HS}; nothing here touches those cases. No minimality is claimed:
\(\dim(P_5)_4\geq254\) is an inequality, we do not determine
\(\dim(P_5)_4\), and we exhibit no generating set for \(P_n\). Oeding's
set-theoretic theorem is untouched --- \(\mathcal V(J_n)=Z_n\) still holds
--- and what fails is the equality of ideals.

The five quartics \(B_1,\dots,B_5\) are not new: they are the five
degree-four invariants of a five-qubit state constructed by Luque and
Thibon \cite{LT}, obtained by fixing one qubit, taking the degree-two
invariant of the remaining quadrilinear form, and then its discriminant.
What is new here is the common value \(12t^{10}\Pi^2\) of their pullbacks
\eqref{eq:pullback}, and with it the membership \eqref{eq:membership}. We
have not found the strict inclusion \(\HD(5)\subsetneq(P_5)_4\), or the
nonradicality of \(J_n\), recorded anywhere: the ideal-theoretic problem is
stated as open in \cite[p.~1194]{HO}, and the literature citing \cite{HS}
and \cite{Oeding} that we were able to examine treats the variety of
principal minors of symmetric matrices set-theoretically. That is an
absence in a search we cannot claim to be exhaustive, and not a proof of
absence.

\begin{thebibliography}{9}

\bibitem{HS}
O.~Holtz and B.~Sturmfels,
\emph{Hyperdeterminantal relations among symmetric principal minors},
J. Algebra \textbf{316} (2007), no.~2, 634--648,
\href{https://doi.org/10.1016/j.jalgebra.2007.01.039}
{doi:10.1016/j.jalgebra.2007.01.039}.

\bibitem{HO}
H.~Huang and L.~Oeding,
\emph{Symmetrization of principal minors and cycle-sums},
Linear Multilinear Algebra \textbf{65} (2017), no.~6, 1194--1219,
\href{https://doi.org/10.1080/03081087.2016.1233932}
{doi:10.1080/03081087.2016.1233932}.

\bibitem{LT}
J.-G.~Luque and J.-Y.~Thibon,
\emph{Algebraic invariants of five qubits},
J. Phys. A \textbf{39} (2006), no.~2, 371--377,
\href{https://doi.org/10.1088/0305-4470/39/2/007}
{doi:10.1088/0305-4470/39/2/007}.

\bibitem{Oeding}
L.~Oeding,
\emph{Set-theoretic defining equations of the variety of principal
minors of symmetric matrices},
Algebra Number Theory \textbf{5} (2011), no.~1, 75--109,
\href{https://doi.org/10.2140/ant.2011.5.75}
{doi:10.2140/ant.2011.5.75}.

\end{thebibliography}

\end{document}